% !PRINTED 383
\weberpage{383}
%% [running head, not body text -- commented out 2026-08-07] Page furniture (running
%% head, folio, signature) is excluded from this edition by standing policy.
%% §\,104. Charaktere und Geschlechter der quadratischen Formen. \hfill 383

Wenn die Klasse $A$ aus der Klasse $A'$, $A''$ komponiert ist,
so ist symbolisch
\[
A = A'A''.
\]

Sei $m'$, $m''$ durch die Klasse $A'$, $A''$ darstellbar, und nehmen
wir $m'$, $m''$ relativ prim, so ist $m = m'm''$ durch $A$ darstellbar
(§ 103), und daraus ergibt sich für die Charaktere die Formel:

\weq{13}{\chi(\delta, A'A'') = \chi(\delta, A')\chi(\delta, A'').}

Nach 1. ist das Symbol $\chi(\delta, A)$ durch irgend eine zu $\delta$
teilerfremde, durch $A$ eigentlich darstellbare Zahl $m$ bestimmt,
und wir definieren daher die Charaktere dieser Zahlen $m$ durch

\weq{14}{\chi(\delta, m) = \chi(\delta, A).}

Es ist aber zweckmäßig, daß man sich noch von der Voraussetzung frei macht, daß $m$ zu $\delta$ relativ prim sei, und dies
kann auf folgende Weise geschehen. An der Voraussetzung, daß
$m$ relativ prim zu $Q$ sei, soll aber festgehalten werden.

Ist $m$ durch die Klasse $A$ eigentlich darstellbar, so können
wir in $A$ einen Repräsentanten

\weq{15}{\varphi = (m, B, C)}

finden. Wir zerlegen $m$ in zwei Faktoren

\weq{16}{m = nn',}

so daß $n$ und $n'$ teilerfremd sind und $n$ relativ prim zu $\delta$, $n'$
relativ prim zu $\delta'$ ist.

Dies ist immer möglich, meist auf mehrere Arten, da nach
Voraussetzung $m$ relativ prim zu $Q$ ist, und also $m$, $\delta$, $\delta'$ keinen
gemeinschaftlichen Teiler haben. Man nehme z. B. in $n$ die in
$\delta'$ aufgehenden Primzahlen in so hohen Potenzen auf, als sie
in $m$ enthalten sind, und außerdem noch die Primzahlpotenzen,
die zu $\delta$ teilerfremd sind.

Dann ist die Form $\varphi$ aus den beiden Formen
\[
(n, B, Cn') \qquad (n', B, Cn)
\]
komponiert, deren Charaktere nach (12) und 4. durch
\[
(\delta, n), \quad (\delta'\ n') % [sic] Gedruckt ohne Komma im zweiten Symbol.
\]
bestimmt sind, und es ist also nach (13):

\weq{17}{\chi(\delta, A) = (\delta, n)(\delta', n').}

Hierin kann aber auch noch die Forderung aufgegeben
werden, daß $m$ durch $A$ eigentlich darstellbar sei; denn ein
gemeinsamer Teiler von $x, y$ (relativ prim zu $Q$) gibt einen
%CONT>>
